\section{Lezione 17} \hfill 09/04/2026\\
pioppelle .... 2 anni diametri tra 4 e 7 cm\\
Non si usano pioppelle più vecchie perché prenderebbero con più difficoltà e perché l'altezza sarebbe eccessiva.\\
Il pioppo non vuole suoli argillosi, eccesso di calcare, ma vuole una buona disponibilità idrica senza avere le radici in falda.\\
La pioppicoltura è normalmente una questione agricola, pertanto i pioppi si piantano in aziende agricole di pianura, normalmente nelle aree meno produttive.\\
Il luogo ideale per la pioppicoltura sono le golene, dove c'è rischio di inondazione: il pioppo sopporta qualche giorno di sommersione.\\
Oggi c'è una polemica in atto sul fatto che i pioppi non andrebbero piantati in golena, poiché queste ultime andrebbero rinaturalizzate.\\
Gli impianti potrebbero essere:
\begin{itemize}
    \item cicli di 2 anni: utilizzando le talee, per la produzione di biomassa da energia. La corteccia brucia malissimo e tende a cementificare gli impianti. Inoltre, dopo due o tre ceduazioni, la ceppaia e la piante muore, non riuscendo a ricacciare. Infine, per creare pioppicoltura da energia vengono usati dei cloni e quindi gli impianti sono molto monotoni: il risultato è il riempimento di animali defogliatori (melasoma populi). L'aspetto è molto simile ad un vivaio di pioppo;
    \item molto usato nel mondo (poco in Italia) il pioppo come tondello da carta, con turni di 5-7 anni. Il diametro obiettivo è di 10 cm?? Il materiale utilizzato sono le talee. Le distanze sono normalmente di 2x3, .... e 3x4;
    \item il metodo tradizionale (italiano) ha un turno tra i 10 e 14 anni, con una tendenza di 10 anni in Friuli e 14 in Piemonte. L'obiettivo sono i toppi da sfoglia (compensati). Il diametro obiettivo è una circonferenza di 100 cm. Con un toppo di 2-2,20 metri.... Le distanze sono circa 35-40 mq a pianta. Il sesto di impianto è definitivo (non sono previsti diradamenti).
    \item il metodo francese ha un diametro obiettivo di 50-60 cm , da sfoglia, con turni di 19-21 anni. Il numero di piante messe a dimora sono sui 100 ad ha. La distanza aumenta perché il pioppo non si sopporta, quando c'è la competizione laterale con altri individui: la distanza quindi aumenta in proporzione al diametro che si vuole ottenere dall'impianto.
\end{itemize}
Il metodo italiano è più agronomico, quello francese è più forestale (perché si tende ad abbandonare). Il metodo francese non è molto apprezzato in Italia, perché il turno è più lungo; inoltre, le chiome tendono a chiudere intorno al decimo anno, lasciando suolo scoperto che non produce. Per chi compra è meglio il metodo francese, perchè in proporzione lo scarto è minore.\\
I pioppi quando arrivano in campo devono essere reidratati (per almeno una settimana): questo facilita l'attecchimento, soprattutto per i cloni caroliniani.\\
I sesti sono il disegno con il quale viene messo giù l'impianto:
\begin{itemize}
    \item il quadrato
    \item il rettangolare: usato in agroforestazione, in modo da usare il suolo scoperto. Non funziona molto, perché il mais prende risorse al pioppo ed il pioppo crea ombra al mais;
    \item il quittonce: il modulo iniziale è un triangolo isoscele. È il più usato.
    \item il settonce: il modulo iniziale è un triangolo equilatero. A parità di piante, occupa meno spazio.
\end{itemize}
L'impianto viene fatto ad inizio primavera- fine autunno, prima che il pioppo entri in stagione vegetativa, con una buca di almeno 0.8-1 metro.\\
L'apparato radicale è avventizio, che tende a fare radici in superficie per l'acqua meteorica e in profondità per l'acqua di falda.\\
Nel primo anno post impianto non si fa nulla: occorre lasciarlo fare radici e chioma. Se necessario, si fanno 1-2 irrigazioni di soccorso. Occorre scegliere bene le stazioni di impianto, limitando così le future irrigazioni.\\
Alla fine del primo anno potrebbero esserci delle fallanze: essendo tutti cloni, dovrebbero essere tutti vivi o tutti morti (ad esclusione di disturbi abiotici, es. trattore). Le fallanze devono essere ripiantate, entro il primo anno.\\
È meglio una pioppella di un anno, perché costa meno e produce più tronco.\\
Dal primo anno in poi, fino a che le chiome non chiudono (4-6 anni), occorre diserbare meccanicamente (discatura): la competizione erbacea potrebbe prendere acqua e risorse. Le discature rischiano di creare ferite alle radici avventizie. Dopo le discature dei primi anni, vengono fatte trinciature.\\
Nei primi anni vengono fatti trattamenti contro gli insetti del legno (saperla, cossus cossus), che generano gallerie e rovinano il legno, e contro altri insetti. Inoltre, si fanno trattamenti in chioma, contro defogliatori e ....\\
Quando chiudono le chiome, vengono fatti trattamenti solo sulla chioma, oppure non si fanno più perché gli atomizzatori diventano efficaci.\\
Arrivati a fine turno si vende il pioppo ed il prezzo viene definito da un collaudatore dell'acquirente. Il prezzo è unitario a quintale, sulla base della qualità del legno:
\begin{itemize}
    \item diametro
    \item assenza di legni di reazione
    \item ferite
    \item difetti, es. nodi
\end{itemize} 
Il prezzo, pre covid, era di 12 euro al quintale di prima qualità e 6 euro di pessima qualità. Il prezzo è ciclico nel tempo.\\
Se si vuole reimpiantare di nuovo un pioppeto, occorre eliminare le vecchie ceppaie. Nell'anno successivo all'utilizzazione, non si dovrebbe piantare nuovamente un impianto, ma sarebbe da aspettare 1-2 stagioni.\\
Un buon sistema per trattare meno i pioppi, di recente concezione, è quello di fare impianti policiclici: sullo stesso appezzamento occorre piantare biomassa da energia (5-7 anni di turno, nocciolo, robinia, platano, carpino), un turno breve di 10-14 anni del pioppo ed un turno lungo di 20 di latifoglie nobili (noce,...). Ogni 3-4 anni si ritraggono prodotti, .... le distanze sono maggiori, le risorse sono maggiori e si fanno molti meno trattamenti (se non azzerati). Il problema è che ci sono poche piante ad ettaro.
\subsection{Noce}
Da sempre, in Italia, è la specie più pregiata.\\
Il noce è una specie che attualmente non si vende, perché è passata di moda.\\
Ha bisogno di almeno 7-8 metri tra una pianta e l'altra: sennò non cresce.\\
Ha un legno diviso in alburno e durame: turni troppo brevi non servono.\\
Richiede un toppo utile di almeno 2 metri (fusto senza rami).\\
Tendenzialmente, nella corteccia rimangono sempre i difetti passati della pianta.  
\subsection{Rimboschimenti}
I termini forestazione e riforestazione sono ambigui. Il termine piantumazione è da giardinieri.\\
I tipi di operazione sono
\begin{itemize}
    \item imboschimento: impianto di alberi su una superficie che non è mai stata bosco, solitamente dove non c'è mai stato suolo (es. campo agrario, frana, cava, discarica);
    \item rimboschimento: ricostruzione di un popolamento che per qualche ragione è andato distrutto (es. disturbi abiotico) o per utilizzazione artificiale;
    \item sottoimpianto: aumento della densità di alberi nel bosco, possono essere della stessa specie che c'è già (velocizzare la successione naturale, es. impianto di abete rosso in un lariceto). Non vuol dire piantare sotto la chioma di altri alberi (la competizione dell'albero già presenta è eccessiva), ma in fianco (negli spazi liberi).
\end{itemize}
Il materiale usato può essere:
\begin{itemize}
    \item il seme: è il migliore, perché la pianta che ne deriva ha un apparato radicale perfetto. Il problema è la grande predazione da parte degli animali. Inoltre, la raccolta dei semi è molto costosa per molte specie arboree. Viene usato per specie con seme molto grosso e pesante. Solo una parte dei semi acquistati sono fertili;
    \item piantine: per specie con semi leggeri, ottenute da seme o da talee. Vengono coltivate in vivaio, selezionando le piantine belle da quelle brutte. Possono essere:
    \begin{enumerate}
        \item a radice nude: molto usato in passato, senza terra. È facile da trasportare, non avendo il peso e volume della terra. Gli apparati radicali vengono protetti. Sono più economici. L'apparato radicale si sviluppa in pieno campo, nelle piante fittonanti è meglio questo metodo perché il fittone può crescere libero. Questo metodo soffre di stress da trapianto, che ci impedisce di piantare in periodo vegetativo (limita nelle operazioni). Presenta un numero elevato di fallanze.
        \item pane di terra (in vaso): molto di moda adesso. Meno o nullo stress da trapianto, che permette di piantare in periodo vegetativo. La terra può essere arricchita con sostanze nutritive. Costa di più, sia per produzione che per trasporto. Non deve avere una crescita costretta, ovvero le radici non devono toccare il vaso (potrebbe soffrire di nanismo o morire). 
        \item materiale moltiplicato in vitro: vengono fatte moltiplicare le gemme in vitro, in modo da non moltiplicare le sue malattie. È molto costoso. Si usa per le specie che non riescono più a moltiplicarsi.
    \end{enumerate}
\end{itemize}
Una piantina per performare al meglio, deve avere un perfetto equilibrio tra la chioma e le radici. Se la chioma è eccessiva, l'apparato radicale piccolo non riesce a nutrire la pianta correttamente. Se la chioma è ....\\
Il meglio per l'equilibrio chioma-radici si ottiene quando la pianta utilizzata per i rimboschimenti è piccola.\\
Il vaso perfetto deve avere:
\begin{itemize}
    \item una forma conica: se la radice tocca il bordo, viene spinta verso il basso;
    \item presenti delle scanalature, in modo che se la radice inizia a spiralare trovi degli impedimenti per scendere 
    \item con degli scoli
\end{itemize}